% Source: https://artofproblemsolving.com/community/c3h2108510_generalization_of_2020_sprint_29
% Stable identity: aops-topic-2108510-q02
Use vectors with the origin at $G$, and write $\vec A,\vec B,\vec C,\vec P$ for the position vectors. Since $G=\dfrac{A+B+C}{3}$ is the origin, $\vec A+\vec B+\vec C=\vec0$.

\begin{asy}
usepackage("amsmath");
size(5.5cm);
defaultpen(fontsize(10pt));
pair A=(0.6,4), B=(0,0), C=(6,0); pair G=(A+B+C)/3; pair P=(4.6,3.3);
draw(A--B--C--cycle,linewidth(1));
draw(A--(B+C)/2,gray+dashed); draw(B--(A+C)/2,gray+dashed); draw(C--(A+B)/2,gray+dashed);
draw(P--A,blue); draw(P--B,blue); draw(P--C,blue);
draw(P--G,red+linewidth(1));
dot(A); dot(B); dot(C); dot(G,red); dot(P,blue);
label("$A$",A,N); label("$B$",B,SW); label("$C$",C,SE); label("$G$",G,SW,red); label("$P$",P,NE,blue);
label("\mbox{}",(min(currentpicture,user=true)+max(currentpicture,user=true))/2,black); // keeps the text after the figure black when LaTeX typesets the labels inline
\end{asy}

\textbf{(a)} For each vertex, $PA^2=\left|\vec A-\vec P\right|^2=\left|\vec A\right|^2-2\vec A\cdot\vec P+\left|\vec P\right|^2$. Adding the three such equations,
\[PA^2+PB^2+PC^2=\left(\left|\vec A\right|^2+\left|\vec B\right|^2+\left|\vec C\right|^2\right)-2\left(\vec A+\vec B+\vec C\right)\cdot\vec P+3\left|\vec P\right|^2=GA^2+GB^2+GC^2+3PG^2 .\]

\textbf{(b)} Let $S=\left|\vec A\right|^2+\left|\vec B\right|^2+\left|\vec C\right|^2=GA^2+GB^2+GC^2$ and $T=\vec A\cdot\vec B+\vec B\cdot\vec C+\vec C\cdot\vec A$. Then
\[0=\left|\vec A+\vec B+\vec C\right|^2=S+2T,\qquad a^2+b^2+c^2=\left|\vec B-\vec C\right|^2+\left|\vec C-\vec A\right|^2+\left|\vec A-\vec B\right|^2=2S-2T .\]
So $T=-\dfrac {S}{2}$ and $a^2+b^2+c^2=2S+S=3S$, which gives $S=\dfrac{a^2+b^2+c^2}{3}$.

\textbf{(c)} By (a), $PA^2+PB^2+PC^2\ge GA^2+GB^2+GC^2$, with equality exactly when $PG=0$, that is, $P=G$. By (b), the smallest value is $\dfrac{a^2+b^2+c^2}{3}$; for sides $5,5,6$ it is $\dfrac{25+25+36}{3}=\boxed{\dfrac{86}{3}}$.
